% Generated draft: requires mathematical and source review.
\documentclass[11pt]{article}
\usepackage[margin=25mm]{geometry}
\usepackage{fontspec}
\setmainfont{Latin Modern Roman}
\usepackage{amsmath,amssymb,mathtools,unicode-math}
\setmathfont{Latin Modern Math}
\usepackage{enumitem,needspace,xurl,hyperref}
\setlength{\parindent}{0pt}
\setlength{\parskip}{5pt}
\setlength{\emergencystretch}{3em}
\newcommand{\N}{\mathbb{N}}
\newcommand{\Z}{\mathbb{Z}}
\newcommand{\Q}{\mathbb{Q}}
\newcommand{\R}{\mathbb{R}}
\newcommand{\C}{\mathbb{C}}
\newcommand{\F}{\mathbb{F}}
\newcommand{\A}{\mathbb{A}}
\newcommand{\PP}{\mathbb{P}}
\newcommand{\E}{\mathbb{E}}
\newcommand{\eps}{\varepsilon}
\newcommand{\dd}{\,\mathrm{d}}
\begin{document}
\section*{UnsolvedMath ID: 800868}
\textbf{Source problem code:} catalogue entry
\subsection{Definitions and mathematical statement}
The supplied source asks to characterize the image of an entropy map \(H\), its intersections with the faces of \(\mathrm{SubMod}_n\), and its fibers. The record does not define the domain of \(H\) or the precise convention for \(\mathrm{SubMod}_n\).

Write \(D\) for the intended domain and \(C=\mathrm{SubMod}_n\) for the intended submodular cone. Once these objects and the map \(H:D\to C\) are specified, the requested problems are to characterize
\[
\operatorname{im}(H)=\{H(p):p\in D\},
\]
to characterize, for each face \(F\) of \(C\),
\[
\operatorname{im}(H)\cap F,
\]
and to describe, for each \(h\in\operatorname{im}(H)\), the fiber
\[
H^{-1}(h)=\{p\in D:H(p)=h\}.
\]
These expressions refer to the actual image; replacing it by its closure would require an additional qualification absent from the supplied statement.

For clarity, a possible finite discrete interpretation would take \(D\) to be a specified family of joint probability laws of \(n\) random variables and set
\[
H(p)_S=-\sum_{x_S}p_S(x_S)\log_b p_S(x_S),
\qquad \varnothing\ne S\subseteq\{1,\ldots,n\},
\]
where \(p_S\) is the marginal law, \(b>1\), and \(0\log_b0=0\), with \(H(p)_\varnothing=0\) if the empty-set coordinate is included. A set function \(h\) is submodular if
\[
h(A)+h(B)\ge h(A\cup B)+h(A\cap B)
\quad\text{for all }A,B\subseteq\{1,\ldots,n\}.
\]
This interpretation is conditional: the supplied record does not specify finite discrete variables, allowed alphabet sizes, the logarithm base, coordinate conventions, or whether \(\mathrm{SubMod}_n\) also imposes normalization and monotonicity. Consequently, it does not determine a unique concrete entropy-map problem.
\subsection{Short English statement}
Characterize the entropy vectors attained by the intended map, determine which attained vectors lie on each face of the intended submodular cone, and describe all inputs producing a given entropy vector. The domain and cone conventions are missing from the supplied excerpt, so the concrete formulation requires review.
\subsection{Sources}
\begin{itemize}
\item Local UnsolvedMath registry: permanent ID \(800868\), source problem code \texttt{catalogue entry}, title \emph{Open Problems in Algebraic Statistics --- General Problem 7}. The supplied record attributes the problem to Sturmfels, with proposed year \(2007\), and labels it open; that status is supplied registry metadata, not an independently verified assessment.
\item Supplied source attribution: Sturmfels, \emph{Open Problems in Algebraic Statistics} (2007), General Problem (\#7), \url{https://arxiv.org/abs/0707.4558}. The background records access on 2026-10-16 and extraction by \texttt{arxiv\_latex}; these are supplied provenance statements, not a claim of independent access or verification.
\item Supplied upstream revision: \texttt{d371ee917ac9647fbb06cd6fc529deec6f3dda82}. The excerpt supplies the request concerning the image, face intersections, and fibers, but omits the definitions of its principal objects.
\end{itemize}
\end{document}
